\section{Auswertung des Jahresabschlusses}
\label{sec:abschluss}

Alle Betr\"age in Millionen Euro, sofern nicht anders angegeben. Grundlage ist die
Mehrjahresreihe 2017--2025 (Tabellen~\ref{tab:ergebnis} und~\ref{tab:zahlung}); das erste
Halbjahr 2026 steht jeweils am Ende.

\ergebnistabelle

\subsection{Umsatz}

\[
\Delta\text{Umsatz} = \frac{\text{Umsatz}_{2025}}{\text{Umsatz}_{2024}} - 1
\]

Der Umsatz fiel 2025 um \UmsatzDelta{} (Tabelle~\ref{tab:ueberblick}). Er liegt damit unter dem
H\"ochststand von \Mio{\UmsatzMax}{EUR} aus dem Jahr \UmsatzMaxJahr. Im ersten Halbjahr 2026
stieg er um \HjUmsatzDelta{} auf \Mio{\HjUmsatz}{EUR}\QL{hjdrei}{HJ}{3}. Der Umsatz folgt bei Neste
vor allem den Produktpreisen und sagt \"uber die Ertragslage wenig: Im zweiten Quartal 2026
erh\"ohten h\"ohere Endproduktpreise den Umsatz um rund \je{\HjPreiseffekt}{Mrd.\,EUR},
geringere Absatzmengen minderten ihn\QR{hjdrei}.

\subsection{F\"uhrender Indikator: Verkaufsmarge von Renewable Products}

Neste ver\"offentlicht keinen Auftragseingang. Die Gr\"o{\ss}e, die das Ergebnis vorwegnimmt,
ist die vergleichbare Verkaufsmarge von Renewable Products je Tonne: Umsatz abz\"uglich
Rohstoff- und Logistikkosten je verkaufter Tonne. Sie lag 2019 bis 2025 zwischen
\je{\MargeXXIV}{USD/t} und \je{\MargeMax}{USD/t} (Median \je{\MargeMedian}{USD/t}),
erreichte im ersten Halbjahr 2026 \je{\HjRpMarge}{USD/t}\Q{HJ}{6} und im zweiten Quartal
\je{\QzRpMarge}{USD/t} nach \je{\QzRpMargeVj}{USD/t} im Vorjahresquartal\Q{HJ}{1}. Das
Halbjahr liegt \"uber jedem Jahr der eigenen Reihe (Rang \Pz{\MargeRangHj}), beim
\MargeFaktorHj-fachen des Medians. Die Verbindung zum Kurs behandelt
Teil~\ref{sec:kurs}.

\subsection{EBITDA-Marge: ausgewiesen und vergleichbar}

Neste f\"uhrt das \emph{vergleichbare} EBITDA als Hauptkennzahl. Es bereinigt das
ausgewiesene EBITDA um Bewertungsgewinne und -verluste der Vorr\"ate, Marktwert\"anderungen
offener Rohstoff- und W\"ahrungsderivate, Ver\"au{\ss}erungsergebnisse und
Versicherungsentsch\"adigungen\Q{GBXXV}{148}.

\begin{table}[htbp]\centering\small
\caption{Ausgewiesenes gegen vergleichbares EBITDA (Mio.~EUR), aus
Tabelle~\ref{tab:ergebnis}.}\label{tab:ebitda}
\begin{tabular}{@{}lrrrl@{}}\toprule
Jahr & EBITDA & Vergl.\ EBITDA & Differenz & h\"oher ist\\\midrule
\tabellenkoerper{data/ebitda_abgleich}
\bottomrule\end{tabular}\end{table}

Das Vorzeichen der Differenz wechselt: In \EbitdaHoeherJahre{} lag das ausgewiesene EBITDA
\"uber dem vergleichbaren, in allen anderen Jahren darunter. Das ist die Spur der
Vorratsbewertung: In Jahren steigender Produktpreise gewinnt der Vorrat an Wert, und das
ausgewiesene Ergebnis enth\"alt diesen Gewinn; in Jahren fallender Preise der umgekehrte
Fall. \bk{Das vergleichbare EBITDA ist deshalb die bessere Gr\"o{\ss}e f\"ur die
Ertragskraft, aber es ist eine Gr\"o{\ss}e des Unternehmens; der Kassenzufluss in
Tabelle~\ref{tab:zahlung} enth\"alt beide Effekte.}

\subsection{Operatives EBITDA}

Dieser Bericht f\"uhrt das vergleichbare EBITDA als operatives EBITDA, weil Neste es seit dem
ersten Quartal 2022 als Hauptkennzahl der Ertragskraft verwendet\Q{GBXXI}{161} und weil nur
diese Gr\"o{\ss}e \"uber die Jahre gleich abgegrenzt ist. Es stieg 2025 um
\EbitdaDeltaVoll{} auf \Mio{\VerglEbitdaVoll}{EUR}; in den letzten zw\"olf Monaten bis Juni 2026
erreichte es \Mio{\EbitdaLtm}{EUR}. Das ist Rang \Pz{\EbitdaRangLtm} der Reihe
\EbitdaJahre: 2022 und 2023 lag es h\"oher.

\subsection{Nettogewinn}

Das Periodenergebnis betrug 2025 \Mio{\ErgebnisVoll}{EUR} nach einem Verlust im Vorjahr, im
ersten Halbjahr 2026 \Mio{\HjNettogewinn}{EUR}\Q{HJ}{3}; das Ergebnis je Aktie des Halbjahrs
war \je{\HjEps}{EUR} nach \je{\HjEpsVj}{EUR}. Auf zw\"olf Monate gerechnet ergibt sich ein
Ergebnis je Aktie von \je{\EpsLtm}{EUR}.

\subsection{Dividende und Dividendenrendite}

F\"ur 2025 wurden \je{\DividendeVorschlag}{EUR} je Aktie ausgesch\"uttet\Q{GBXXV}{238},
wie im Vorjahr; das ist eine Aussch\"uttungsquote von \Pz{\AusschuettungsQuoteNeste} des
Ergebnisses je Aktie\Q{GBXXV}{146}. Die Dividendenrendite betr\"agt zum Jahresendkurs 2025
\DivRenditeJahresende{}, zum heutigen Kurs \DivRenditeHeute. Die h\"ochste Dividende der Reihe
war \je{\DpsMax}{EUR} f\"ur \DpsMaxJahr, einschlie{\ss}lich einer Sonderdividende.

In den f\"unf Jahren 2021--2025 flossen \Mio{\CfoSumme}{EUR} aus dem Gesch\"aft zu. Dem
standen \Mio{\InvestSumme}{EUR} Investitionen, \Mio{\LeasingSummeFuenf}{EUR} Leasingtilgung
und \Mio{\DivSummeFuenf}{EUR} Dividenden gegen\"uber. Die L\"ucke von
\Mio{\FinanzierungsLuecke}{EUR} wurde mit Schulden geschlossen; die Nettofinanzschulden
stiegen von 2021 bis 2024 um \Mio{\NettoschuldenAnstieg}{EUR}.
\bk{Die Dividenden der Jahre 2022 und 2023 sind zu einem erheblichen Teil mit Fremdkapital
bezahlt worden. Die K\"urzung auf \je{\DividendeVorschlag}{EUR} war die Folge, nicht die
Ursache der Schw\"ache.}

\subsection{Aktienkursverlauf}

\kursabbildung

Der Kurs erreichte \je{\KursMax}{EUR} im Jahr \KursMaxJahr{} und fiel bis auf
\je{\KursMin}{EUR} im Jahr \KursMinJahr{} (Kennzahlenseiten). Der heutige Kurs liegt
\KursVomHoch{} unter dem H\"ochststand.

\subsection*{Kennzahlen\"uberblick}

\ueberblicktabelle
\zahlungstabelle

Der freie Zufluss schwankte 2017--2025 zwischen \Mio{\FcfMin}{EUR} und \Mio{\FcfMax}{EUR};
in \FcfNegativJahre{} von neun Jahren war er negativ, jeweils in einem Jahr gro{\ss}er
Investitionen. Sein Median liegt bei \Mio{\FcfMedian}{EUR}. Das Jahr 2025 liegt auf Rang
\Pz{\FcfRangVoll} der eigenen Reihe, die letzten zw\"olf Monate mit \Mio{\FcfLtm}{EUR} auf
Rang \Pz{\FcfRangLtm}. In diesem LTM-Wert steckt ein Aufbau des Umlaufverm\"ogens um
\Mio{\HjNwcAufbau}{EUR} im ersten Halbjahr 2026, vor allem Vorr\"ate f\"ur die Stillst\"ande
der zweiten Jahresh\"alfte\Q{HJ}{2}; ohne die Bewegung des Umlaufverm\"ogens w\"aren es
\Mio{\FcfLtmOhne}{EUR}.

Ohne Wachstumsinvestitionen, also nur nach Instandhaltung und Leasing, h\"atte das Gesch\"aft
in den Jahren \ErnteJahre{} im Median \Mio{\ErnteMedian}{EUR} abgeworfen. Der Unterschied zum
Median des freien Zuflusses ist der Preis des Wachstums: Singapur und Rotterdam.
